\documentclass[../DoAn.tex]{subfiles}
\begin{document}
\section{Đặt vấn đề}
\label{section:1.1}
Ngày nay, cùng với sự phát triển mạnh mẽ của khoa học máy tính và trí tuệ nhân tạo, \gls{its} đang đóng vai trò quan trọng trong việc quản lý và điều hành giao thông đô thị. Trong đó, \gls{alpr} là một trong những thành phần quan trọng nhất, được ứng dụng rộng rãi trong giám sát an ninh, xử lý vi phạm giao thông, thu phí tự động và quản lý bãi đỗ xe.

Phần lớn các hệ thống ALPR thương mại hiện nay được thiết kế để hoạt động trong môi trường có kiểm soát (constrained environments), nơi camera được lắp đặt ở vị trí cố định với góc quay chuẩn, ánh sáng đồng đều và phương tiện di chuyển ở tốc độ chậm hoặc dừng hẳn. Tuy nhiên, khi triển khai vào môi trường thực tế không ràng buộc (unconstrained environments), chẳng hạn camera giám sát trên đường cao tốc hay camera hành trình, hiệu năng của các hệ thống này suy giảm nghiêm trọng. Nguyên nhân chính là dữ liệu thu được thường chứa đựng nhiều yếu tố gây nhiễu: biến đổi ánh sáng, thời tiết xấu, khoảng cách xa khiến độ phân giải biển số thấp, rung lắc camera và đặc biệt là hiện tượng mờ nhòe do chuyển động (motion blur) khi phương tiện di chuyển với tốc độ cao.

\begin{figure}[H]
    \centering
    \includegraphics[width=0.9\textwidth]{Hinhve/challenges_grid.jpg}
    \caption{Một số thách thức thường gặp đối với ảnh cắt biển số trong thực tế (từ trái qua phải, từ trên xuống dưới: mờ nhòe do chuyển động, lóa sáng do đèn pha và thiếu sáng/nhiễu hạt)}
    \label{fig:challenges}
\end{figure}
Từ những thách thức đó, việc xây dựng một hệ thống có khả năng nhận dạng chính xác biển số xe ngay cả khi phương tiện đang chuyển động trong điều kiện ngoại cảnh phức tạp là một yêu cầu cấp thiết. Xuất phát từ nhu cầu thực tiễn này, bài toán ``Phát hiện và nhận dạng biển số xe chuyển động trong các nguồn video'' được lựa chọn làm đề tài nghiên cứu và phát triển trong đồ án.

\section{Các giải pháp hiện tại và hạn chế}
\label{section:1.2}
Các giải pháp ALPR thế hệ đầu chủ yếu dựa vào kỹ thuật xử lý ảnh truyền thống như phát hiện biên, phân tích đường viền, hoặc sử dụng các đặc trưng thủ công như HOG và Haar. Nhóm phương pháp này phụ thuộc lớn vào các ngưỡng cấu hình cứng, thiếu tính khái quát hóa và nhạy cảm với sự thay đổi của ngoại cảnh.

Sự bùng nổ của học sâu (Deep Learning), tiêu biểu là các mạng nơ-ron tích chập (CNN), mô hình phát hiện đối tượng họ YOLO và các mạng nhận dạng ký tự quang học (OCR), đã khắc phục phần lớn các nhược điểm trên. Tuy nhiên, các giải pháp hiện hành thường tiếp cận bài toán dưới góc độ xử lý ảnh tĩnh đơn khung hình (single-image ALPR). Khi áp dụng trực tiếp vào luồng video, việc ép buộc mô hình đưa ra dự đoán trên một khung hình đơn lẻ bị suy giảm chất lượng mà bỏ qua ngữ cảnh thời gian xung quanh thường dẫn đến kết quả nhận dạng sai lệch.

Một số giải pháp gần đây đã bắt đầu tích hợp thuật toán theo dõi đa đối tượng để bám vết phương tiện, nhưng vẫn tồn tại hai hạn chế chính. Thứ nhất, thuật toán bám vết truyền thống dễ mất dấu hoặc nhảy định danh (ID switch) khi phương tiện thay đổi vận tốc đột ngột hoặc bị che khuất. Thứ hai, phần lớn các giải pháp chưa thiết lập được cơ chế đánh giá chất lượng khung hình để chắt lọc thông tin dư thừa theo thời gian (temporal redundancy); việc nhận dạng lặp lại trên mọi khung hình mà thiếu bộ lọc chất lượng không chỉ lãng phí tài nguyên mà còn khiến kết quả đầu ra thiếu ổn định.

\section{Mục tiêu và định hướng giải pháp}
\label{section:1.3}
Nhằm khắc phục các hạn chế nêu trên, đồ án đề ra ba mục tiêu cụ thể:

\begin{enumerate}
    \item \textbf{Xây dựng và chuẩn hóa các bộ dữ liệu đa nhiệm cho biển số Việt Nam.} Hệ thống dữ liệu được phân rã thành các tập chuyên biệt phục vụ từng bài toán lõi, bao gồm: 40.465 ảnh gán nhãn \gls{obb} cho phát hiện biển số và hơn 33.000 ảnh cắt biển số cho huấn luyện OCR. Các tập dữ liệu được thiết kế để bao phủ đa dạng điều kiện ngoại cảnh phức tạp trong thực tế.

    \item \textbf{Đề xuất luồng xử lý nhận dạng đa khung hình dựa trên học sâu.} Hệ thống phát hiện và bám vết quỹ đạo phương tiện, sau đó định vị và cắt vùng biển số bên trong từng track. Một bộ định tuyến chất lượng (Quality Router) tự động phân loại khung hình đủ rõ để OCR ngay và khung hình cần tích lũy thêm bằng chứng. Kết quả cuối cùng được tổng hợp từ nhiều khung hình thông qua thuật toán đối sánh chuỗi thời gian ký tự (Character Time-series Matching).

    \item \textbf{Triển khai thành hệ thống phần mềm end-to-end trên nền tảng Web.} Ứng dụng cho phép tải lên video, kết nối camera giám sát qua luồng RTSP, quan sát tiến trình nhận dạng và tra cứu lịch sử kết quả kèm hình ảnh chứng cứ.
\end{enumerate}

\section{Đóng góp của đồ án}
\label{section:1.4}
Các đóng góp chính của đồ án gồm: (i) xây dựng và chuẩn hóa các tập dữ liệu phục vụ nhiều tác vụ trong hệ thống ALPR, bao gồm dữ liệu phát hiện biển số dạng OBB, dữ liệu OCR biển số Việt Nam và dữ liệu đánh giá chất lượng khung hình; (ii) thiết kế pipeline nhận dạng biển số trên video theo hướng đa khung hình, kết hợp phát hiện phương tiện, phát hiện biển số, tracking, lọc chất lượng, OCR và bỏ phiếu kết quả theo track; (iii) huấn luyện và tích hợp các mô hình thành phần phù hợp với bối cảnh biển số Việt Nam, trong đó có mô hình phát hiện biển số YOLOv8-OBB, mô hình phân loại chất lượng khung hình và mô hình OCR SmallLPR-Line-CTC; (iv) xây dựng hệ thống Web end-to-end cho phép tải video, kết nối nguồn RTSP, theo dõi tiến trình xử lý, xem kết quả nhận dạng theo phương tiện và lưu lại lịch sử nhận dạng kèm ảnh chứng cứ; và (v) đánh giá thực nghiệm các mô hình thành phần và pipeline tổng thể, từ đó chỉ ra hiệu quả của việc khai thác thông tin dư thừa theo thời gian trong bài toán nhận dạng biển số xe trên video.

Kết quả đạt được của đồ án là một hệ thống thử nghiệm hoàn chỉnh, có khả năng tiếp nhận và xử lý video giao thông thực tế, hiển thị kết quả nhận dạng sau quá trình xử lý và lưu trữ dữ liệu phục vụ tra cứu lại. Phần đóng góp trọng tâm không nằm ở một mô hình đơn lẻ, mà ở việc tổ chức các mô hình phát hiện, theo dõi, OCR và hậu xử lý thành một luồng xử lý thống nhất cho bài toán nhận dạng biển số xe Việt Nam trên video.

\section{Bố cục đồ án}
\label{section:1.5}
Báo cáo đồ án được tổ chức thành sáu chương như sau.

Chương 1 giới thiệu đề tài, trình bày bối cảnh, các hạn chế của hệ thống ALPR hiện tại, mục tiêu nghiên cứu và các đóng góp chính.

Chương 2 trình bày nền tảng lý thuyết phục vụ cho hệ thống. Nội dung chương bao gồm bài toán nhận dạng biển số xe trên video, các mô hình phát hiện đối tượng, kiến trúc YOLOv5 và YOLOv8-OBB, thuật toán theo dõi đa đối tượng BoT-SORT, các hướng tiếp cận OCR.

Chương 3 mô tả phương pháp đề xuất. Chương này trình bày tổng quan pipeline xử lý, quá trình thu thập và xử lý dữ liệu, cách huấn luyện các mô hình thành phần gồm phát hiện phương tiện, phát hiện biển số, truy vết phương tiện, phân loại chất lượng khung hình và giới thiệu mô hình OCR SmallLPR-Line-CTC. Phần cuối chương mô tả các bước hậu xử lý, xác thực định dạng biển số và bỏ phiếu kết quả theo nhiều khung hình.

Chương 4 trình bày quá trình xây dựng và triển khai hệ thống. Nội dung chương bao gồm yêu cầu hệ thống, sơ đồ use case, kiến trúc tổng quan, các công nghệ sử dụng, thiết kế frontend, backend và API, cơ sở dữ liệu, kiểm thử chức năng và cấu hình triển khai.

Chương 5 đánh giá thực nghiệm các mô hình và pipeline tổng thể. Chương này mô tả môi trường thử nghiệm, tiêu chí đánh giá, kết quả của các mô hình thành phần như mô hình phát hiện phương tiện, mô hình phát hiện biển số và mô hình OCR, sau đó phân tích hiệu quả của pipeline trên các kịch bản video và các trường hợp lỗi thường gặp.

Chương 6 tổng kết các kết quả đã đạt được, nêu lại những đóng góp chính của đồ án và trình bày các hướng phát triển tiếp theo nhằm cải thiện độ chính xác, tốc độ xử lý và khả năng triển khai của hệ thống trong môi trường thực tế.
\end{document}